%------------------------------------------------------------------------------%
\begin{question}
  Tente calcular a inversa da matriz
  \quad
  \(
    A =
    \begin{pmatrix}
      1 & 2 \\
      2 & 4
    \end{pmatrix}
  \)
\end{question}

%------------------------------------------------------------------------------%
\begin{resolution}{Matriz Aumentada}
  Matriz aumentada \([\,A\,|\,I\,]\)
  \pause
  \[
    \begin{amatrix}[2]{2}
      1 & 2 & 1 & 0 \\
      2 & 4 & 0 & 1
    \end{amatrix}
  \]

  \pause
  Queremos transformar a parte da esquerda na identidade $I_2$
\end{resolution}

%------------------------------------------------------------------------------%
\begin{resolution}{Escalonamento}
\begin{columns}
\column{0.3\linewidth}
  \[\;\;\;
    \begin{amatrix}[2]{2}
      1 & 2 & 1 & 0 \\
      2 & 4 & 0 & 1
    \end{amatrix}
  \]
  \pause
\column{0.7\linewidth}
  Eliminamos os elementos abaixo do primeiro pivô
  \\ \pause
  \(
    L_2 \leftarrow L_2 - 2 L_1
  \)
\end{columns}
\vspace{-\baselineskip}
\[
  \pause
  L'_2
  \;=\;
  L_2 - 2 L_1
  \pause
  \;=\;
  \begin{amatrix}[2]{2}
    2 & 4 & 0 & 1
  \end{amatrix}
  - 2
  \begin{amatrix}[2]{2}
    1 & 2 & 1 & 0
  \end{amatrix}
  \pause
  \;=\;
  \begin{amatrix}[2]{2}
    0 & 0 & -2 & 1
  \end{amatrix}
\]
\pause
\[
  \begin{amatrix}[2]{2}
    1 & 2 & 1 & 0 \\
    0 & 0 & -2 & 1
  \end{amatrix}
\]
\end{resolution}

%------------------------------------------------------------------------------%
\begin{resolution}{Solução}
  \[
    \begin{amatrix}[2]{2}
      1 & 2 & 1 & 0 \\
      0 & 0 & -2 & 1
    \end{amatrix}
  \]

  \pause
  O escalonamento da matriz $A$ produziu uma linha nula

  \pause
  A matriz $A$ não tem posto completo

  \pause
  A matriz $A$ é singular

  \pause
  A matriz $A$ não possui inversa
\end{resolution}

%------------------------------------------------------------------------------%
\endinput
